\section{Plane curves and the finite maps}\label{sec:plane-cover}

We first construct two finite maps to a common plane.  Their pulled-back
curves will supply isometric integral Seifert forms and the monodromy
sectors needed later.  We also retain the homology maps induced by each
projection, because those maps determine the integral boundary condition
in the fourfold construction.

Set
\begin{align*}
 h(X,T)&=T(X-T^2)(X-2T^2),\\
 q_1(u,w)&=(u-w-w^8)(u-w-2w^8)(u-w^2),\\
 q_2(u,w)&=q_1(u,w)-u^2 .
\end{align*}
We use two graph maps to the \((X,T)\)-plane:
\begin{align*}
 \pi_0(u,w)&=\bigl(\varphi(u,w),u\bigr),
    &\varphi(u,w)&=u^2+q_1(u,w),\\
 \pi_1(x,w)&=\bigl(x,\psi(x,w)\bigr),
    &\psi(x,w)&=u(u-w^2)q_2(u,w),\quad u=x+w.
\end{align*}
Write \(d_i=h\circ\pi_i\). For \(g\in\{h,d_0,d_1\}\), let \(F_g\)
denote the chosen sufficiently small compact oriented Milnor page of
\(g\). Set \(\mathcal H=F_h\), \(\mathcal D_i=F_{d_i}\), and put
\(H=H_1(\mathcal H;\mathbb Z)\), \(D_i=H_1(\mathcal D_i;\mathbb Z)\).
Throughout, \(s(a,b)\) is the ordered intersection of a filling of
\(a\) on a reference page with a filling of the positive neighboring-page
transport of \(b\), in the acyclic total Milnor tube.

\begin{lemma}[Plane data]\label{lem:plane-data}
Each \(\pi_i\) is a finite map germ of degree four. The curves \(d_i=0\)
are reduced, have isolated singularities at the origin, and have the same
six-branch embedded type. Their integral Milnor Seifert forms are
isometric, and each has Milnor number \(51\).
Let \(\Phi_5(t)=t^4+t^3+t^2+t+1\).  If \(M_h\) is the monodromy on
\(H_{\mathbb Q}=H\otimes\mathbb Q\), then
\[
 \dim_{\mathbb Q}\ker(M_h-I)=2,\qquad
 \dim_{\mathbb Q}\ker\Phi_5(M_h)=4.
\]
The monodromy on either \(D_i\otimes\mathbb C\) is semisimple at
eigenvalue \(1\) and at every primitive fifth root of unity.
\end{lemma}

\begin{proof}
\noindent\emph{Finite maps and branch contacts.}
At \(u=0\),
\[
 \varphi(0,w)=-w^4(1+w^7)(1+2w^7),
\]
and at \(x=0\),
\[
 \psi(0,w)=-w^4(1-w)+O(w^{19}).
\]
Weierstrass preparation in \(w\) therefore makes both graph projections
finite of degree four.

In the first graph \(d_0=u q_1q_2\). Introduce invertible local
coordinates \(x=u-w,\ t=u\). Two branches of \(q_1=0\) satisfy
\[
 x=A_e(t)=e\,t^8+O(t^{15}),\qquad e=1,2.
\]
The branch \(u=w^2\) satisfies \(t=x^2+O(x^3)\), and \(u=0\) gives \(t=0\).
The other two local branches arise from \(q_2=0\). Substituting \(u=vw^2\)
gives
\[
 q_2(vw^2,w)=w^4(v-1-v^2)+O(w^5).
\]
The two roots of \(v^2-v+1=0\) are simple; hence the implicit-function
theorem gives two branches
\[
 t=B_f(x)=f\,x^2+O(x^3),\qquad
 f=\frac{1+i\sqrt3}{2},\ \frac{1-i\sqrt3}{2}.
\]
These exhaust the local branches of \(q_2=0\): indeed
\(q_2(u,0)=u^2(u-1)\), so Weierstrass preparation in \(u\) has degree two.
Together with \(f=0,1\), these are four branches \(t=B_f(x)\).
All six are smooth and distinct.

In the second graph,
\[
 d_1=\psi(x,w)\bigl(x-\psi(x,w)^2\bigr)
                    \bigl(x-2\psi(x,w)^2\bigr).
\]
The four branches of \(\psi=0\) are exactly the four \(t=B_f(x)\)
above. Since \(\psi(0,w)=-w^4+O(w^5)\), the last two factors each have
a unique smooth branch \(x=e w^8+O(w^9)\), \(e=1,2\), and these read
\(x=e t^8+\text{higher terms}\) in \(t=x+w\). Thus both \(d_i\) are units
times products of two \(x=A_e(t)\) factors and four \(t=B_f(x)\) factors
with the same first distinguishing coefficients.

The two \(A\)-branches meet with multiplicity \(8\); any pair of
\(B\)-branches meets with multiplicity \(2\); a branch from each group
meets transversely. Hence
\[
 \delta(d_i)=8+\binom42\,2+2\cdot4=28,\qquad
 \mu(d_i)=2\delta(d_i)-6+1=51.
\]
\smallskip
\noindent\emph{The integral plane-form comparison.}
Interpolate only the terms above the displayed leading coefficients
and interpolate the nonvanishing units through units. This joins both
germs to the same leading product
\[
 d_\circ=(x-t^8)(x-2t^8)
       \prod_{f\in\{0,1,(1+i\sqrt3)/2,(1-i\sqrt3)/2\}}(t-fx^2).
\]
We now verify uniform isolation and boundary transversality along the
interpolation, which yield an oriented isotopy of its Milnor tubes.
In the sector \(|x|\le 2|t|\), if \(x/t^8\) stays bounded, weighted
rescaling has the noncritical model
\(t^4(x-t^8)(x-2t^8)\) at \(t=1\); its factors have distinct simple
roots, and its weighted Euler derivative excludes other critical
points. If \(|x/t^8|\to\infty\), all \(t-B_f(x)\) factors are
\(t(1+o(1))\), both \(x-A_e(t)\) factors are \(x(1+o(1))\), and
\(x\,\partial_x d/d\to2\). In the sector \(|t|\le2|x|\), the analogous
rescaling for bounded \(t/x^2\) gives
\(x^2\prod_f(t-fx^2)\); for \(|t/x^2|\to\infty\) one has
\(t\,\partial_t d/d\to4\). These estimates and their first
derivatives are uniform in the interpolation parameter after the
ball is shrunk. The six branch parametrizations also meet small
spheres transversely and remain disjoint there. The associated
Milnor tubes and boundary open books consequently form an oriented
isotopy. Transporting cycles and filling chains through it preserves
the integral Seifert pairing. Thus \(D_0\) and \(D_1\) have isometric
integral Seifert forms.

\smallskip
\noindent\emph{The required monodromy sectors.}
The plane-form comparison is now established. We finish by computing
the two sectors used in the arithmetic argument and excluding Jordan
blocks there. The three smooth branches of \(h\) have pairwise intersections
\(1,1,2\); thus \(\mu(h)=6\). The weighted rotation
\((X,T)\mapsto(\zeta^2X,\zeta T)\), \(\zeta^5=1\), preserves \(h\)
and has order five. Its nontrivial generator has no fixed point on
\(F_h\). Lefschetz therefore gives trace \(1\) on \(H_{\mathbb Q}\). Over
\(\mathbb Q\), a representation of the cyclic group of order five
is a sum of trivial and four-dimensional cyclotomic representations.
Its rank and trace equations force ranks \(2\) and \(4\), respectively.

It remains to justify the asserted Jordan condition for \(d_\circ\).
On a sufficiently small sphere, write its argument in the sector
\(|x|\le2|t|\) as the argument of
\[
 t^4(x-t^8)(x-2t^8)
 \prod_f(1-fx^2/t),
\]
and in \(|t|\le2|x|\) as the argument of
\[
 x^2\prod_f(t-fx^2)
 \prod_{e=1,2}(1-et^8/x).
\]
The displayed correction products are uniformly close to \(1\)
where used. Smooth cutoffs in the ratio \(|x|/|t|\) remove their
small arguments by an isotopy of argument fibrations. To see
regularity through that isotopy, use the angular fields of the
weighted rotations with rates \((8,1)/20\) on the \(t\)-side and
\((1,2)/10\) on the \(x\)-side. Each differentiates its corresponding
leading phase at rate one. On their common sector, both differentiate
\(2\arg x+4\arg t\) at rate one; the derivatives of the small
correction terms can be made arbitrarily small. Their convex angular
interpolation therefore remains positively transverse to every
intermediate phase. Near the binding only the appropriate weighted
field is used and is tangent to the binding. Thus the isotopy
preserves the open-book monodromy.

This periodic-piece description follows the plane-link calculus of \cite[Sections~11 and~13]{EisenbudNeumann1985}. The final phase in the overlap is \(2\theta_x+4\theta_t\).
Each regular phase level there has two annuli, both parametrized
by the radial ratio and the free angle \(\theta_t\), with the same
orientation. A single phase circuit exchanges their cores
\(\gamma_1,\gamma_2\) while preserving their angular orientations.
Outside the overlap, twenty circuits of either weighted rotation
return identically. Across each overlap annulus the
\(\theta_t\)-return varies from one turn on the \(t\)-side to four
turns on the \(x\)-side; hence the twentieth return makes the same
signed threefold Dehn twist on both annuli. The oriented boundary
of the \(t\)-side subsurface is \(\gamma_1+\gamma_2\) modulo actual
page-boundary circles. The latter lie in the radical of the surface
intersection form. Write \(M_\circ\) for the monodromy of
\(d_\circ\). Then
\[
 \operatorname{im}(M_\circ^{20}-I)\subseteq
 \mathbb Q(\gamma_1-\gamma_2),\qquad
 M_\circ(\gamma_1-\gamma_2)=-(\gamma_1-\gamma_2).
\]
Since \(M_\circ^{20}-I\) commutes with \(M_\circ\), it vanishes on the generalized
eigenspaces at \(1\) and at primitive fifth roots. The polynomial
\(z^{20}-1\) has simple roots; \(M_\circ\) is therefore semisimple on
those sectors.
\end{proof}

The two projections retain different degree profiles on the branches.
Here each tuple lists the degrees of the branches above the indicated
component of \(h=0\):
\[
\begin{array}{c|c|c}
\text{base branch} & \pi_0 & \pi_1\\ \hline
T=0     &(4)&(1,1,1,1)\\
X=T^2   &(1,1,2)&(4)\\
X=2T^2  &(2,2)&(4)
\end{array}
\]
We use the integral isometry of \(D_0\) and \(D_1\) only as an isometry
of their abstract forms; no compatibility with the two projections to
\(H\) is assumed.  The next lemma records the two maps on homology
associated with each projection.

\begin{lemma}[Finite-cover adjunction]\label{lem:cover-adjunction}
For \(i=0,1\), the page maps induce
\(\pi_{i*}:D_i\to H\) and a full-preimage transfer
\(\pi_i^!:H\to D_i\) satisfying
\[
 \pi_{i*}\pi_i^!=4\,\mathrm{id}_H .
\]
If \(s_h,s_{d_i}\) are the Seifert forms computed by ordered
filling-chain intersection in the total Milnor tubes, then
\[
 s_{d_i}(\pi_i^!a,b)=s_h(a,\pi_{i*}b),\qquad
 s_{d_i}(b,\pi_i^!a)=s_h(\pi_{i*}b,a).
\]
Consequently, after tensoring with \(\mathbb Q\),
\[
 D_{i,\mathbb Q}=(\ker\pi_{i*})_{\mathbb Q}
       \ \perp\ \pi_i^!H_{\mathbb Q},\qquad
 \operatorname{pr}_i(l)=l-\tfrac14\pi_i^!\pi_{i*}l
\]
is the projection onto \((\ker\pi_{i*})_{\mathbb Q}\) for both orders
of the form.
\end{lemma}

\begin{proof}
For an analytic germ \(f\) that is not locally constant, an ordinary
Milnor tube \(B_\epsilon\cap f^{-1}(D_\delta)\), with
\(0<\delta\ll\epsilon\ll1\), has contractible total; this is the
consequence of Milnor's proof recorded in
\cite[Theorem~1.8 and Remark~1.9]{Massey2015LeCycles}.
We compare the finite graph cuts used here with those ordinary tubes.
Choose finite proper graph representatives over one sufficiently small
closed base ball, and choose ordinary Milnor balls for \(h,d_i\) inside
the corresponding cuts. The branch parametrizations above allow regular
radii to be chosen so that the zero set in each closed intervening region
is exactly a disjoint union of branch annuli. The cut boundaries are
smooth near these annuli, and each annulus is transverse to both
boundaries. The function is a complex submersion normal to each annulus.
A field along the annulus directed from the inner boundary to the outer
one therefore extends, by this submersion, to a field tangent to the
nearby function levels. The disjoint flows give a value-preserving
collar of a neighborhood of the entire collection of annuli.
Properness makes the closed intervening region compact. Outside a
neighborhood of these annuli the absolute value of the function has a
positive lower bound, so a sufficiently small value disk confines the
intervening tube to that collar. Attaching the collar to the ordinary
tube preserves its homotopy type and gives regular lateral boundaries.
Thus these selected cut totals are acyclic. Choose nearby pages whose
base branch values lie in their interiors.
Represent a base class by loops avoiding these finitely many
values; the full preimage is a cycle upstairs.  This operation is
Poincar\'e--Lefschetz dual to pullback in relative first cohomology of
the pages, so it depends only on the homology class of the loops.
Counting the four sheets gives \(\pi_{i*}\pi_i^!=4\).

For the first pairing identity, fill the base loop \(a\) on the reference
page in the acyclic base total tube, and fill the positive neighboring-page
transport of \(b\) upstairs. Pull the first oriented surface chain back
through the finite holomorphic map, and push the second filling forward.
Put both fillings
and their boundary collars in general position. The base filling
meets the generic branch curve at isolated points, and the other
filling can be moved off their preimages; hence intersections relevant
to the pairing occur away from ramification. At each such intersection
the usual oriented projection formula counts the same local number
before and after pushforward. The finite map is over the value disk,
so it also commutes with positive nearby-page transport. This proves
the first identity; reversing the two filling arguments proves the
second. Finally, \(\pi_i^!H\) is rationally nondegenerate, since its
restricted form is \(4s_h\) by adjunction. Its two-sided orthogonal
complement is \(\ker\pi_{i*}\), and the displayed formula is the
corresponding projection.
\end{proof}

We have obtained the plane-form isometry together with a separate
pushforward, transfer, and rational projection for each graph.  The next
section uses the graph equations to construct the two smooth fourfold
germs and to read their ordinary orders in analytic coordinates.
